% !TEX program = xelatex
\documentclass{英语试卷}
\begin{document}
\examtitle{盐城市、南京市 2024-2025 学年度第一学期期末调研测试}{英语}
\par 盐城市、南京市 2024-2025 学年度第一学期期末调研测试\par
\par 英语试题卷\par
\begin{notice}
\par 1. 答卷前，考生务必将自己的姓名、准考证号填写在答题卡上。\par
\par 2. 回答选择题时，选出每小题答案后，用铅笔把答题卡上对应题目的答案标号涂黑。如需改动，用橡皮擦干净后，再选涂其他答案标号。回答非选择题时，将答案写在答题卡上，写在本试卷上无效。\par
\par 3. 考试结束后，将本试卷和答题卡一并交回。\par
\end{notice}
\parttitle{第一部分 听力（共两节，满分 30 分）}
\par 做题时，先将答案标在试卷上，录音内容结束后，你将有两分钟的时间将试卷上的答案转涂到答题卡上。\par
\sectitle{第一节（共5小题：每小题1.5分，满分7.5分）}
\par 听下面 5 段对话。每段对话后有一个小题，从题中所给的 A、B、C 三个选项中选出最佳选项。听完每段对话后，你都有 10 秒钟的时间来回答有关小题和阅读下一小题。每段对话仅读一遍。\par
\question{1. What will the man probably do this afternoon?}
\begin{choices}(3)
  \item Have some coffee.
  \item Prepare for a test.
  \item Take an exam.
\end{choices}
\question{2. What is the man doing?}
\begin{choices}(3)
  \item Walking through a forest.
  \item Visiting the woman's house.
  \item Planting flowers in his garden,
\end{choices}
\question{3. Where does the conversation probably take place?}
\begin{choices}(3)
  \item At a police station.
  \item By the roadside.
  \item At an animal shelter.
\end{choices}
\par 4 What does the man decide to do regarding the new laptop? A. Buy one right now. B. Wait for a better deal. C. Anticipate the release.\par
\question{5. What are the speakers talking about? 微信公众号 IAI English}
\begin{choices}(3)
  \item A lecture.
  \item Some countries.
  \item Traditional customs
\end{choices}
\sectitle{第二节（共15小题：每小题1.5分，满分22.5分）}
\par 听下面 5 段对话或独白。每段对话或独白后有几个小题，从题中所给的 A、B、C 三个选项中选出最佳选项。听每段对话或独白前，你将有时间阅读各个小题，每小题 5 秒钟；听完后，各小题将给出 5 秒钟的作答时间。每段对话或独白读两遍。 听第 6 段材料，回答第 6、7 题。\par
\question{6. What will the woman probably do first?}
\begin{choices}(3)
  \item Go to a modern museum.
  \item Tour a historic building.
  \item Visit a memorial park
\end{choices}
\question{7. What is the relationship between the speakers?}
\begin{choices}(3)
  \item Strangers.
  \item Travel companions.
  \item Tour guide and visitor.
\end{choices}
\par 听第 7 段材料，回答第 8 至 10 题\par
\question{8. What mistake did the woman make?}
\begin{choices}(3)
  \item She forgot to send the desks. 微信公众号 IAI English
  \item She failed to send the desks in time.
  \item She didn't send the other kind of desks
\end{choices}
\question{9. What docs HDC request the speakers to do now?}
\begin{choices}(3)
  \item Give the money back.
  \item Replace all the desks.
  \item Apologize to them
\end{choices}
\question{10. What will the woman do to solve the problem?}
\begin{choices}(3)
  \item Give the customer a call.
  \item Visit the customer instantly.
  \item Arrange the repair of desks
\end{choices}
\par 听第 8 段材料，回答第 11 至 13 题，\par
\par 1 1. How many people will stay at the hotel? A. 6. B. 11. C. 12.\par
\question{12. When will the man check in?}
\begin{choices}(3)
  \item In early November.
  \item In mid-November.
  \item In late November.
\end{choices}
\question{13. What will the man do next?}
\begin{choices}(3)
  \item Confirm the check-out date.
  \item Ask about the course sessions
  \item Offer his personal information 听第 9 段材料，回答第 14 至 17 题。
\end{choices}
\question{14. What did the woman think of the film?}
\begin{choices}(3)
  \item It was exciting.
  \item It was boring
  \item It was scary.
\end{choices}
\question{15. What is the man going to do this evening?}
\begin{choices}(3)
  \item Watch a film.
  \item Visit a friend.
  \item Go bowling
\end{choices}
\question{16. Where is the modern art museum?}
\begin{choices}(3)
  \item Next to Riverside Bowling
  \item Across from the man's flat.
  \item On the north side of the river.
\end{choices}
\question{17. When will the man probably leave his flat?}
\begin{choices}(3)
  \item At 6:45 p.m.
  \item At 7:00 p.m.
  \item At 7:15 p.m.
\end{choices}
\par 听第 10 段材料，回答第 18 至 20 题。\par
\question{18. How do smart homes make life easier for some people?}
\begin{choices}(3)
  \item By increasing Internet speed.
  \item By reducing the cost of living.
  \item By automating household tasks
\end{choices}
\question{19. Why are people concerned about smart homes?}
\begin{choices}(3)
  \item They may affect privacy.
  \item They are difficult to set up.
  \item They may be too expensive
\end{choices}
\question{20. What will the speaker discuss in the next session?}
\begin{choices}(3)
  \item Concrete operations.
  \item Detailed pros and cons.
  \item Future development.
\end{choices}
\parttitle{第二部分 阅读理解（共两节，满分 50 分）}
\sectitle{第一节（共15小题：每小题2.5分，满分37.5分）}
\par 阅读下列短文，从每题所给的 A、B、C 和 D 四个选项中，选出最佳答案，\par
\passagelabel{A}
\par Thank you for your interest in The UCLA Health Volunteer Summer Program. Once you are accepted into the program, you may begin your volunteering in the following critical areas:\par
\par Wayfinding\par
\par Wayfinders are stationed at hospital entrances and greet each visitor/patient along with providing wheelchair assistance. Wayfinders must be friendly and have good communication skills 2-hour shifts are available Monday-Friday between 8:00 a.m. —4:00 p.m. Due to the standing nature of this position, 4 hours will be rewarded for 2 hours of Wayfinding service.\par
\par Floaters\par
\par Floaters help deliver patient mail and newspapers. Assignments can change each shift depending on needs of the department, Volunteers are assigned to the Volunteer Office for one 4-hour shift per week Monday-Friday between 9:00 a.m. —4:00 p.m.\par
\par Patient Transport\par
\par Patient Transport volunteers are responsible for discharging patients and helping the patient transport staff. This is a great opportunity for students to see all areas of the hospital and get clinical exposure. Shifts are available seven days a week anytime between 8:00 a.m. — 8:00 p.m.\par
\par Applying to the Summer Volunteen Program\par
\par Step 1: Sign up for an Information Session on zoom, Sign-ups are on a first come, first serve basis.\par
\par Step 2: Submit all required application forms online for review and approval.\par
\par Step 3: Attend one In-Person Orientation (培训), There is only one session, and we will not have make-up sessions.\par
\par Step 4: Complete 50 volunteer hours by August 31, 2025.\par
\question{21. What can we learn about Wayfinding?}
\begin{choices}(2)
  \item Time bonus will be given.
  \item It is available on weekends.
  \item Assignments change with needs.
  \item It trains volunteers to find ways.
\end{choices}
\question{22. What is an advantage of Patient Transport?}
\begin{choices}(2)
  \item It offers clinical training.
  \item It assists in treating patients
  \item It provides flexible shifts.
  \item It helps show patients around.
\end{choices}
\question{23. What are applicants expected to do when applying for the program?}
\begin{choices}(2)
  \item Sign up as soon as possible.
  \item Track volunteer hours
  \item Attend a make-up orientation.
  \item Review application forms
\end{choices}
\passagelabel{B}
\par More than 20 years ago, Yung-Chi Cheng, a Yale professor in drug development for cancer, had a pioneering idea: what if he could unlock the potential of ancient Chinese medicines for treating cancer? What if he could design botanical drugs that would make traditional cancer treatments work better? No one had done it before. Fellow researchers and experts advised him to change course, because developing botanical drugs was too complicated and too risky. But Cheng didn't let the idea go.\par
\par Now, in a landmark moment in cancer research, Cheng and research partners are launching the first international clinical trial for a botanical drug, YIV-906. Shwu-Huey Liu, a postdoctoral researcher in Cheng's Yale lab, spent hours leafing through ancient Chinese texts at Yale's library seeking the properties of Chinese herbs. She returned with Huang Qin Tang, a 1,800-year-old treatment for stomach illnesses, based on which the team developed a drug and tested it on one thousand mice.\par
\par What followed were years of additional testing of the drug's effectiveness for a range of cancers in multiple human studies involving over 200 patients. The positive effects displayed again and again, YIV-906 not only reduced the side effects of radiation treatment, but also led to faster recovery and longer survival rates.\par
\par All they need now are larger data sets. When the results of the international trial are available about three years from now, YIV-906 may find itself next in a long line of breakthrough drugs developed from natural products, which includes a malaria (疟疾) treatment developed from the Chinese herb ginghao, the discovery of which led to a Nobel Prize for chemist Tu Youyou.\par
\par It has been a long journey, but if Cheng and his team are successful, it could create a new approach for treating cancer. They called the new approach "WE" medicine, a combination of Western medicine and Eastern medicine. "It's a totally new model. I've met with a lot of doubt, but think the results will speak for themselves," Cheng said.\par
\question{24. What did fellow researchers initially think of Cheng's research?}
\begin{choices}(2)
  \item Creative.
  \item Comprehensive.
  \item Conventional
  \item Challenging.
\end{choices}
\question{25. Why did Liu refer to ancient Chinese texts?}
\begin{choices}(3)
  \item She wanted to question ancient treatment methods
  \item She tried to document the history of Chinese herbs 3. She desired to find inspiration for treating diseases
  \item She determined to find a cure for stomach illnesses
\end{choices}
\question{26. What does paragraph 3 mainly talk about?}
\begin{choices}(2)
  \item Steps of the drug testing.
  \item The application of YIV-906.
  \item Effects of the clinical trial.
  \item The effectiveness of YIY-906
\end{choices}
\question{27. Why does the author mention Chinese herb ginghao?}
\begin{choices}(2)
  \item To illustrate the value of botanical drugs
  \item To advocate the adoption of "WE" medicine.
  \item To praise Tu's contributions to drug development
  \item To highlight its success in medicine development.
\end{choices}
\passagelabel{C}
\par Media literacy (素养) is finally getting some late-arriving attention. Whether it's the increasing awareness of political disinformation or the threat of Al, discussions about media literacy are on the rise. The adoption of media literacy in the classroom, however, stagnates.\par
\par The driving force for providing kids with media literacy skills often gets lost in the moral panic over technology and kids. Fear has overrun reason, and people can't decide whether to blame social media, mobile devices, tech companies, or the Internet. They are concerned with claims that technology and social media are the causes of teenage mental health issues despite the weak associations. In response, too many solutions are "ostrich policies (鸵鸟政策)” that want to ban devices and shut down social media but don't address more fundamental issues\par
\par Policies that set age limits on social media access and restrict devices during school to minimize disturbance can be helpful but are incomplete solutions. Social media didn't invent bullies (霸凌) or the enormous effect of peer pressure. Social media platforms may be a new battleground, but they are not the source. Unfortunately, many current policies ignore the need to teach kids how to manage these challenges. There is frighteningly little attention paid to preparing kids to make good decisions when they will, unavoidably, get online.\par
\par On the bright side, the rapid development of AI and increasing attention to the amount of misinformation online have helped normalize the acceptance of media literacy as a core skill. While suggestions to limit inappropriate device use, increase face-to-face socialization, and engage in outdoor play arc beneficial, they do not build the core skills that will keep kids safe offline and are fundamental to well-being and life success: good citizenship, self-awareness and self-control. Those things take teaching, whether you are online or outside.\par
\par Media literacy is an essential life skill that, once learned, can be applied to any content. While waiting for public education to catch up, parents and educators can build media literacy skills at home with numerous resources, ranging from family technology contracts and discussion guides to easy activities that can be done at home with kids.\par
\question{28. What does the underlined word "stagnates" in the first paragraph probably mean?}
\begin{choices}(2)
  \item Gets far ahead.
  \item Falls apart suddenly.
  \item Falls well short.
  \item Gets started officially.
\end{choices}
\question{29. What causes people's moral panic over technology and kids?}
\begin{choices}(2)
  \item Belief in the rapid growth of new technologies
  \item Concern over cutting kids off from the Internet
  \item Inability to balance the use of social media and health.
  \item Fear of technology causing mental health issues in teenagers.
\end{choices}
\question{30. What does paragraph 3 of the text focus on?}
\begin{choices}(2)
  \item The sources of peer pressure.
  \item The benefits of minimizing disturbance.
  \item The problems of current policies.
  \item The disadvantages of social media platforms
\end{choices}
\question{31. Which of the following belongs to media literacy?}
\begin{choices}(2)
  \item Reducing social media exposure.
  \item Applying critical thinking to media.
  \item Increasing face-to-face interactions.
  \item Recognizing the dangers of social media
\end{choices}
\passagelabel{D}
\par Decision fatigue (决策疲劳) is a psychological phenomenon where the quality of decisions declines after a long session of decision-making. It is caused by the limited mental energy we have for making choices each day. While many of us believe we are logical beings capable of weighing every option carefully, decision fatigue shows how our mental resources can be used up over time.\par
\par One of the key studies on this phenomenon was conducted in the context of parole hearings (假释听证会). Researchers found that judges were more likely to approve of parole early in the day or after a meal break, As the day progressed and the judges made more decisions, they were increasingly likely to maintain their previous decisions, denying parole simply because it required ess mental effort.\par
\par Decision fatigue not only affects major life choices but also everyday decisions. For instance, people often report struggling to decide what to eat for dinner after a long day. Similarly, too many choices, such as when looking through streaming platforms or online stores, can leave individuals feeling confused, leading them to abandon the process entirely or settle for a less satisfying option.\par
\par To overcome decision fatigue, some well-known figures have adopted strategies to simplify their daily routines. For example, Facebook founder Mark Zuckerberg is famous for wearing the same style of gray T-shirts to minimize unnecessary choices. Similarly, Nobel Prize-winning physicist Albert Einstein was known for owing multiple identical suits to avoid wasting energy on deciding what to wear each day. In a different context, author Haruki Murakami has shared that he maintains a strict daily routine, including specific writing hours and exercise, to free up mental capacity for his creative work.\par
\par Understanding decision fatigue can help us structure our lives more effectively. Prioritizing important decisions earlier in the day or setting up routines to minimize unnecessary choices are strategies to preserve mental energy. By managing our mental resources wisely, we can make sound decisions and reduce the stress associated with overwhelming choices.\par
\question{32. What leads to decision fatigue according to the text?}
\begin{choices}(2)
  \item The lack of time to make decisions.
  \item The difficulty of making logical choices
  \item The stress from making important decisions.
  \item The limited mental energy for decision-making.
\end{choices}
\question{33. What does the study on parole hearings reveal?}
\begin{choices}(2)
  \item Judges are much stricter after long breaks.
  \item Decision fatigue can lead to unfair outcomes.
  \item Fatigue causes judges to avoid complicated cases.
  \item External factors have no impact on decision-making.
\end{choices}
\question{34. What does the author suggest readers shopping online do?}
\begin{choices}(2)
  \item Make quick decisions.
  \item Compare different products.
  \item Limit the number of choices
  \item Choose the most popular products.
\end{choices}
\question{35. What is the purpose of adopting routines?}
\begin{choices}(2)
  \item To reduce the stress of daily work.
  \item To increase decision-making quality.
  \item To enhance overall life satisfaction.
  \item To create more decision-making chances 阅读下面短文，从短文后的选项中选出能填入空白处的最佳选项。选项中有两项为多余选项，
\end{choices}
\par Achieving Success in Journalistic News Writing\par
\par No book, article or blog post can turn an ambitious writer into a successful journalist. All that it can do is to help the writer learn to write. The primary skill in learning to write is practice. \blanknum{36}\par
\par Among the most convenient subjects are the events and news in the world around us. Such topics are what skilled journalists write about, and that is why we call this form of writing "journalistic writing".\par
\par Journalistic writing is not a special literary gymnastics that requires a special talent and an expansive vocabulary. \blanknum{37} It is using language to tell what is going on in the world. As such, it requires no special talent beyond seeing and understanding what is taking place around us, to grasp its significance and relation to other events, and to tell others what we have seen and heard.\par
\par \blanknum{38} One of the chief characteristics is timeliness and relevance. It has little time to talk about the past, except in comparison with the present. It has little time to foretell the future, except in pointing out the significance of the present. Its subject is today. It is primarily concerned with facts. The writer considers himself a recorder of the things other persons are doing and thinking.\par
\par Another basic part of journalistic writing is that the writer prepares it for a definite purpose. Some writers, like novelists, write to use beautiful words and sentences. \blanknum{39} However, the writer who favors a journalistic style uses English simply to tell what he has seen and heard; he realizes that his reader cares more for the things he tells than for the way in which he tells and expresses them.\par
\par Success in journalistic writing demands two important skills. One is the inborn skill to see and understand. The other is the skill to tell what the writer has seen in a way that makes the reader see the same thing. \blanknum{40} If the writer is eager to excel, he must expand his knowledge and education so he can better understand and evaluate the things he sees and hears.\par
\begin{choices}(1)
  \item It will not do them any good to write news items
  \item To these imaginative literary writers, they favor expression.
  \item To practice writing, one must have something to write about.
  \item An ambitious writer can develop both skills if he lacks them.
  \item The writer must write it now or never—without giving up quality
  \item Journalistic writing differs mainly in its subject matter and purpose.
  \item Instead, it is the plainest, most reserved form of English composition.
\end{choices}
\parttitle{第三部分 语言知识运用（共两节，满分 30 分）}
\sectitle{第一节（共15小题；每小题1分，满分15分）}
\par 阅读下面短文，从短文后各题所给的 A、B、C 和 D 四个选项中，选出可以填入空白处的最佳选项，\par
\par When Mysha decided to improve her fitness in late 2019, taking up running was a no-brainer. She lives in a beautiful town on Scotland's east coast, so she found \blanknum{41} wasn't so bad, with breathtaking scenery.\par
\par Yet even though Mysha quickly fell in love with running, something started to \blanknum{42} her—the beaches and paths she ran along were covered in litter. \blanknum{43} she was always doing a little litter-picking, but when she saw an American runner online talking about plogging, a new \blanknum{44} emerging among runners, she was excited to try it out for herself, Plogging \blanknum{45} in Sweden, and the word comes from the Swedish word ploka, meaning "to \blanknum{46} while jogging". Mysha started following a(n) \blanknum{47} online called Environmental Fitness and began to take an shopping bag with her on runs and fill it with \blanknum{48}\par
\par It wasn't long until she uncovered some of the upsetting \blanknum{49} of littering. She started to find \blanknum{50} chemicals contained in batteries leaking into the environment and small animals \blanknum{51} in plastic bottles struggling to chew their way out. These discoveries made Mysha even more \blanknum{52} about litter picking. She joined her local litter-picking group and soon was asked to \blanknum{53} leadership of the group. Despite the big volumes of litter being dropped, Mysha and her fellow litter-pickers \blanknum{54} to get out and plog.\par
\par "The thought of keeping other people's spirits up and seeing their \blanknum{55} spread to more people keeps me going. Together, we can achieve amazing things!" Mysha says.\par
\cz{41}{jogging}{skating}{cycling}{hiking}
\cz{42}{puzzle}{excite}{bother}{amuse}
\cz{43}{Proudly}{Actually}{Secretly}{Luckily}
\cz{44}{trend}{skill}{lesson}{trick}
\cz{45}{existed}{renewed}{advanced}{originated}
\cz{46}{relax}{pick}{think}{recycle}
\cz{47}{order}{rule}{initiative}{tradition}
\cz{48}{food}{clothes}{tools}{waste}
\cz{49}{truths}{processes}{instances}{consequences}
\cz{50}{rare}{poisonous}{organic}{unstable}
\cz{51}{raised}{trapped}{placed}{buried}
\cz{52}{curious}{worried}{passionate}{confused}
\cz{53}{take over}{insist on}{turn down}{rise to}
\cz{54}{ceased}{hesitated}{attempted}{continued}
\cz{55}{actions}{dreams}{hobbies}{memories}
\sectitle{第二节（共10小题；每小题1.5分，满分15分）}
\par 阅读下面短文，在空白处填入 1 个适当的单词或括号内单词的正确形式。\par
\par As the wheel of time turns, traditional practices are witnessing a comeback, Gua sha, an ancient Chinese healing technique, has been gaining \blanknum{} popular) in recent years as a holistic approach to wellness and beauty, because it offers various benefits to modern wellness seekers.\par
\par On TikTok, a before-and-after video of Gua sha has been viewed more than 11.8 million times. Meanwhile, influencers have started to share their Gua sha routines on some other media apps. Attracting \blanknum{57} (grow)attention from global audiences. Moreover, Western beauty stores and \blanknum{58} (celebrity) have jumped at the fashion and started favoring Gua sha or selling its tools.\par
\par Factually, Gua sha has been practiced in China since the Ving Dynasty and is known to have muscle pain relieved \blanknum{59} nprove blood circulation. Originally, it \blanknum{60} (use) on the body, like arms, legs, and back; however, nowadays, it's commonly used to improve the \blanknum{61} (face) skin health, The practice is performed using a smooth-edged tool to gently scrape (刮) areas of the body, \blanknum{62}, according to Traditional Chinese Medicine (TCM), can encourage \blanknum{63} healthy flow of energy and blood.\par
\par Gua sha, a TCM practice, also \blanknum{64} (serve) as a cultural practice, now has \blanknum{65} (incredible) found its way to the West due to its potential health benefits. Undoubtedly, promoting TCM practices could be a much-needed opportunity for others to learn more about Chinese cultures.\par
\parttitle{第四部分 写作（共两节，满分 40 分）}
\sectitle{第一节（满分15分）}
\par 假定你是李华，上周六你所在的城市举办了马拉松赛事，你参加了其中的迷你马拉松比赛。请你给美国朋友 Alec 写一封邮件分享这次经历，内容包括：\par
\par (1) 参赛经过;\par
\par (2) 你的收获。\par
\par 注意:\par
\par (1) 写作词数应为 80 左右;\par
\par (2) 请按如下格式在答题纸的相应位置作答。\par
\par Dear Alec,\par
\par I'm writing to share with you a mini-marathon I participated in last Saturday.\par
\par Yours sincerely, Li Hua\par
\sectitle{第二节（满分25分）}
\par 阅读下面材料，根据其内容和所给段落开头语续写两段，使之构成一篇完整的短文。\par
\par I had always been cautious around Max, He was the quiet kid in our grade, the one who sat at the back of the classroom with his hood (连衣帽) pulled low over his face. Max rarely spoke to anyone, and when he did, his words were rudely brief and often misunderstood. Rumors (谣言) about him spread like wildfire—some said he got into fights after school, others claimed he had once stolen a teacher's wallet, I didn't know if any of them was true, but I stayed away, just like everyone else.\par
\par Then one Friday morning, Max didn't show up to class. A low voice of curiosity spread through the room, and soon enough, theories about his absence started flying. "He got suspended," whispered a girl next to me. "I heard he struck a guy outside the library." I wasn't surprised, and nodded in agreement.\par
\par Later that day, I found myself at the park, enjoying the first warm day of spring. That's when I saw him—Max—sitting alone on the swings. His usual guarded expression was gone, and replaced by something I hadn't expected: sadness. Against my better judgment, I walked over.\par
\par "Hey," I said awkwardly.\par
\par Max looked up, surprised. "What do you want?"\par
\par "I... just saw you here and thought I'd say hi."\par
\par He gave me a strange look but said nothing. I was about to walk away when I noticed a wrinkled piece of paper in his hand, "What's that?" I asked.\par
\par For a moment, he hesitated. Then, with a shrug, he handed it to me. It was a hospital discharge form, and the name on it wasn't his it was his younger sister's. I looked at him, confused.\par
\par "She' been sick for months," Max explained, his voice barely heard. "Yesterday was bad. My mom couldn't leave work, so I stayed at the hospital and looked after her all night."\par
\par It hit me then how wrong we had all been about Max. The tough outward appearance, the distance—he wasn't trying to push people away. He was carrying a weight none of us understood. 注意:\par
\par (1) 续写词数应为 150 左右;\par
\par (2) 请按如下格式在答题卡的相应位置作答，\par
\writingpara{Paragraph 1: The next Monday, Max returned, but rumors remained.}
\writinglines{8}
\end{document}
